% !TeX root = ../papers/nc-p3.tex
\section{特異な種数一曲線上の残余点と固有直線入射の貼合}
\label{sec:singular-residual-gluing}
\begin{theorem}\label{thm:singular-residual-gluing}
特異な八余因子核 $E$ の有限鉛筆を考える。$P=L\eta$、$A_u=H^0(P I_u)$ として
$A$ は $E$ 上階数3・次数 $-4$ の束である。入力の階数1スキーム全体は $M=\mathbb P_E(A)$ 上の
階数3束の切断の零スキーム $Y$ に同型である。$H$ を入力超平面束とする。
切断の束 $F$ は数値的に $[F]=4[\mathcal O(H)]-[\mathcal O]$ と同じChern指標を持つ。
特異点上や比例する二切断の上でも同型は冪零を含めて成立する。
\end{theorem}
\begin{proof}[Cartier四点環における乗法]
$A$ は $H^0(P)\otimes\mathcal O_E\to P$ の評価核である。
任意の非零 $a\in H^0(P)$ の零因子 $D_a$ は、特異点を含む場合も長さ4のCartier因子である。
$L|_{D_a}$ への評価は同型である。非自明な次数零線束 $L/P$ の全コホモロジーが零だからである。
入力パラメータ上のこの因子も有限平坦である。有限非分岐の場合 $Q=L\eta^{-1}$ と置く。
$D_a$ の台で消えない $b_*$ を $H^0(Q)$ に取ると $N(a,b_*)$ は可逆であり、$N(a,b)$ は
これに $D_a$ 上の $b/b_*$ の乗法を合成したものである（評価双対側の規約なら転置）。
一般の単純根で二変数核を対角に評価すればこの等号が得られ、普遍有限平坦因子上の束写像の等号として
$b_*$ が単元の全開集合に延びる。可逆性は二項結果式、または一般根での行列式の恒等式による。

節点二捩れでは $a_\varepsilon=a+\varepsilon b$ が $P_\varepsilon$ の切断の一次持上げである。
$D_a$ への局所成分 $b_{\rm loc}$ は自明化の一次変更で $a$ の倍しか変わらない。
上の乗法表示で使う比は、同じ $a$ の参照持上げに対するこの成分の比となる。
ここでも可逆な参照は全ての $a$ で取れる。実際参照を $P$ の切断だけ変えて各台点での零を避けられる。
その行列式は滑らかな共通零点を印した結果式である。全持上げの空間ではこの入射の像は既約因子、
$a$ を一般に固定すれば独立な四根の四超平面が出る。より直接、$a_0\ne0$ の一般四根の対角公式で、
評価行列式のVandermonde因子を除くと $\det N(a,b)$ は二項結果式 $\operatorname{Res}(a,b)$ の定数倍である。
両者は全次数8である。滑らかな共通根を印した既約因子はこの結果式全体である。両端での共通消滅は
$a_0=b_0=0$ の余次元2であり、一般の滑らかな共通根は単純で他に共通根を持たないからである。
特異点を含む $D_a$ でも $a_0=0$ なら $b_4=-b_0$ なので $b_{\rm loc}(p)\ne0$ は両端での非消滅となる。
従って可逆な参照を用い、単純根の微分対角公式を同じ普遍束で延長できる。

任意の局所Artin底 $S$ で、階数1なら乗法の像は有限自由四点環内の階数1直和因子である。
像を割って得る共通零スキーム $Z$ は $S$ 上平坦長さ3、像への環作用は残余点 $u:S\to E$ を定める。
$E_S$ 上 $I_u=(I_{D_a}:I_Z)$ であり、分数イデアルとして
\[
                           I_Z^\vee\simeq P I_u .
\]
実際局所方程式 $a$ を使えば右のcolonは $a\,\operatorname{Hom}(I_Z,\mathcal O)$ である。
非分岐で $b$ を先に取っても同じ $I_Z=(a,b)$ となり、$Q I_v$ とも同型である。
この比較は $I_Z^\vee$ の標準切断 $1$ に対応する $a,b$ を互いに移す。
\end{proof}

\begin{proof}[自由でない点イデアルを含むAbel写像]
種数一Gorenstein整曲線上の次数1の階数一ねじれなし層 $D_1$ は $h^1=0,h^0=1$ を持つ。
$D_1\to\mathcal O$ の非零射は単射となり次数に反するためである。非零切断の余核は長さ1であり、
双対を取れば $D_1$ は一点イデアルの双対である。これはArtin底上の平坦族にも、底の線束倍を除いて通用する。
実際コホモロジーと底変換により切断は一本の束、繊維での単射性により余核は有限平坦長さ1である。
相対Gorenstein双対でその第一Extは同じグラフ上の線束となる。ねじれなし層の双対は底変換と両立し
二重双対で元に戻る。これは各繊維で最大Cohen--Macaulay層の正のExtが消えることから、
底の平方零拡大に沿い自由分解を双対化しても従う。従って各次数零線束の捩りは、
この次数1表示を介して特異点を含む $E$ 上の並進を定める。

非分岐での $v$ は $u$ のこの並進であり、$u$ 全体で $PI_u$ と $QI_v$ の同定は底の線束 $J$ 倍である。
同定の任意性はスカラーに限られる。繊維で固定した次数1層の唯一の切断を保つ自己同型は恒等であり、
上記の正規化した族の比較が底の線束だけを与えるからである。正規化上では既出の残余一次因子を
入れ替える同定 $j:A_u\to A'_v$ が全点で取れ、従って $\deg J=0$ である。
$M$ 上で $Ta$ を $j(a)$ の生成する直線で割った三成分が求める切断である。
逆にこれが零なら二切断からの $I_Z^\vee$ の切断がスカラー単元倍で同じになる。具体的には
$PI_u$ 内の $a$ に対する $\mathcal O\to PI_u$ の双対の像が $I_Z$ である。
余核の長さは3で平坦（各繊維で非零切断が単射）であり、$b$ からも同じ三点が出る。
乗法の階数は高々1、かつ元の行列の十成分に基点はないので正確に1となる。

節点二捩れでは $I_Z=(a,b_{\rm loc})$ を一定に一次変形すると $a_\varepsilon$ が $Z_\varepsilon$ 上で消える。
$D_{a_\varepsilon}$ の中で残余を取ると長さ1の平坦な点 $u_\varepsilon$ となる。これはCartier四点環の
相対完全跡で平坦な三点商の直交補が直線になるためである。
先と同じ式 $P_\varepsilon I_{u_\varepsilon}\simeq (P I_u)_\varepsilon$ は、$u_\varepsilon$ が所定の一次並進であることを示す。
逆に恒等に特殊化するこの二層の同定を選ぶと、定数切断 $a$ を移して一次持上げ $j(a)$ を得る。
同定の違いは $1+\varepsilon\lambda$ のスカラーなので、微分の差は $\lambda a$ である。
従って指定持上げと $j(a)$ の差を $a$ の直線で割れば大域的な三成分となる。
零なら定数の $Z_\varepsilon$ を持ち、先の乗法表示で階数1となる。
どちらも $Y$ から入力への射で残余点とスカラーが一意に逆構成されるので、
全ての局所Artin底で一致し、スキームとして同型である。
非分岐の三束は $H$ でひねった四つの自明成分から次数零の底の線束を除いたもの、
分岐の三束では除くのは自明線束であり、数値的主張を得る。
\end{proof}

\begin{proposition}\label{prop:singular-horizontal-gluing}
正規化上のスカラー入射で底を支配する曲線の一般重複度の和を $j$ とする。
非分岐節点では $j\le1$、ただし $t^2=-1$ だけは $j=2$ が可能候補で、そのとき二定数は互いに負となる。
尖点と節点二捩れでは $j\le1$ である。$j=1$ の水平部分は滑らかな射影直線で、
$j=2$ の水平部分は二つの射影直線を導手の二点で節点状に貼ったものである。
正規化スカラーを固定すると各水平曲線は入力の直線または二次曲線である。
\end{proposition}
\begin{proof}
正規化の $u$ 上の三空間同士の同定を用いると、スカラー無限遠では包含が単射なので水平成分の
スカラーは定数である。一般核が二次元なら各 $u$ で入力直線が同じスカラー平面を覆ってしまう。
実際同一直線が固定ならその全四次式が全 $u$ で消える。従って一般核は一次元である。
水平曲線の入力平面性と次数上界は定理\ref{thm:deficient-plucker-centers}の台の議論による。

導手で $A_p$ は $a_0=0$ の三空間である。$T$ と比較同定との差を相手三空間の外へ射影する写像は
ここで非零である（$a_0,b_0$ の独立性）。分岐でも $j$ の微分を引いた後のこの条件は $b_0=0$ となる。
従って $p$ の近傍でその核は階数2束 $K$ である。$j$ の同定と局所分裂を用い
$K$ への射影まで取れば、入射のスカラーは二行自己準同型の固有値となる。
その特性多項式は底上の正則なモニック二次式である。これは単元で一行消去した残りの
三行二列鉛筆の一つの最大行列式なので、水平根の一般重複度込みの積で割れる。

節点の両端で同じ底の自明化に戻した水平根は $\lambda,t^2\lambda$、節点二捩れでは
$\lambda,\lambda+1$ となる。正規化表示の $r=0,s=0$ における $a,b$ を比較すればよい。
両値は異なるので一つの水平根でも導手の固有値は二つとも単純である。
$j=2$ なら両端で二根の集合が同じでなければならず、非零二根を保つ非恒等の倍率は $-1$ のみ、
非零平行移動では不可能である。尖点の局所正規化変数を $r$、$s=1$ とすれば比較同定の比は
$(r-(1+2tr)z)/(r-z)=1+2tr+O(r^2)$ である。
実際の底に正則な自明化は $r$ の一次係数を持たない。$j=2$ なら二根の積は非零で一次係数も非零となり、
底の特性多項式の定数係数になれない。$j=1$ ならこの非零一次項により水平枝は尖点上でも滑らかに入る。
節点では二端のスカラーが異なるため単一水平枝は自己貼合しない。二水平枝の場合は二端が交差して貼られ、
底の節点への二枝の射が各枝で局所同型であるから交叉長は各1である。
導手以外では各水平枝は基底の正規化上の切断であり、異なる定数同士は交わらない。
\end{proof}

\path{computations/verify_singular_residual_gluing.py} は節点二型の端点のスカラーと、尖点での同定およびその二乗の一次項を検算する。
